\chapter{บทนำสู่สถานะที่สี่ของสสาร}
\label{chap:ch01}

% =========================================================================
% แผนบริหารการสอนประจำบทที่ 1 (Modular Lesson Plan Boxes - 2-Page Standard)
% =========================================================================

\section*{แผนบริหารการสอนประจำบทที่ 1}
\addcontentsline{toc}{section}{แผนบริหารการสอนประจำบทที่ 1}

% --- Page 1 of Lesson Plan ---
\begin{planbox}{หัวข้อเนื้อหาประจำบท}
\begin{enumerate}
    \item สถานะที่สี่ของสสารและการเปลี่ยนสถานะจากของแข็งสู่พลาสมา
    \item องค์ประกอบของพลาสมาและกลุ่มอนุภาคว่องไวทางเคมี
    \item นิยามและคุณสมบัติพื้นฐานทางฟิสิกส์ของพลาสมา
    \item การจำแนกประเภทพลาสมาตามสมดุลทางอุณหพลศาสตร์
    \item บทบาทและความสำคัญของพลาสมาในงานเกษตรกรรมและสิ่งแวดล้อม
\end{enumerate}
\end{planbox}

\begin{planbox}{วัตถุประสงค์เชิงพฤติกรรม}
\begin{enumerate}
    \item อธิบายกลไกการเปลี่ยนสถานะของสสารจนเกิดเป็นสถานะพลาสมาได้อย่างถูกต้อง
    \item ระบุองค์ประกอบสำคัญ 6 ชนิดภายในพลาสมาความดันบรรยากาศได้อย่างครบถ้วน
    \item คำนวณค่าความยาวเดอบาย ระดับการแตกตัวเป็นไอออน และแปลงหน่วยพลังงานอิเล็กตรอนโวลต์ได้อย่างถูกต้อง
    \item จำแนกความแตกต่างระหว่างพลาสมาสมดุลความร้อนและพลาสมาไม่สมดุลความร้อนได้อย่างชัดเจน
    \item วิเคราะห์ข้อได้เปรียบของการประยุกต์ใช้พลาสมาเย็นในงานเกษตรและสิ่งแวดล้อมได้อย่างเป็นระบบ
\end{enumerate}
\end{planbox}

\newpage

% --- Page 2 of Lesson Plan ---
\begin{planbox}{วิธีสอนและกิจกรรมการเรียนการสอน}
\begin{enumerate}
    \item บรรยายเชิงปฏิสัมพันธ์ประกอบสื่อมัลติมีเดียและแบบจำลองทฤษฎีการกำบังเดอบาย
    \item อภิปรายกลุ่มย่อยเกี่ยวกับการเปรียบเทียบข้อดีและข้อจำกัดของพลาสมาอุณหภูมิห้อง
    \item สาธิตการสร้างพลาสมาเจ็ทและการคายประจุแบบกั้นด้วยไดอิเล็กทริกจำลอง
    \item ศึกษาค้นคว้ากรณีศึกษาการใช้น้ำกระตุ้นด้วยพลาสมาเพื่อทดแทนปุ๋ยเคมีไนโตรเจน
    \item สรุปบทเรียนร่วมกันและทำแบบฝึกหัดทบทวนท้ายบท
\end{enumerate}
\end{planbox}

\begin{planbox}{สื่อการเรียนการสอน}
\begin{enumerate}
    \item เอกสารประกอบการสอนและตำราวิชาการ พลาสมาฟิสิกส์เพื่อการเกษตรและสิ่งแวดล้อม
    \item สไลด์การสอนอิเล็กทรอนิกส์และภาพประกอบเวกเตอร์ความละเอียดสูง
    \item วิดีโอทัศน์การทดลองการคายประจุพลาสมาความดันบรรยากาศและการวัดสเปกตรัมแสง
    \item ชุดคำถามทบทวนประจำบทและใบกิจกรรมวิเคราะห์บทความวิจัยสากล
\end{enumerate}
\end{planbox}

\begin{planbox}{การวัดและประเมินผล}
\begin{enumerate}
    \item ประเมินจากการมีส่วนร่วมในการอภิปรายและตอบคำถามในชั้นเรียน ร้อยละ 10
    \item ประเมินจากความถูกต้องในการทำแบบฝึกหัดทบทวนท้ายบท ร้อยละ 10
    \item ประเมินจากรายงานการวิเคราะห์กรณีศึกษาเทคโนโลยีพลาสมา ร้อยละ 10
    \item ประเมินจากการสอบวัดผลสัมฤทธิ์ประจำภาคการศึกษา ร้อยละ 70
\end{enumerate}
\end{planbox}

\clearpage

% =========================================================================
% เนื้อหาบทเรียน (เริ่มต้นหน้า 3 อย่างเป็นระเบียบ)
% =========================================================================

\section{สถานะที่สี่ของสสารและการเปลี่ยนสถานะ}

สสารในธรรมชาติที่เราคุ้นเคยในชีวิตประจำวันประกอบด้วยสามสถานะหลัก ได้แก่ ของแข็ง (Solid) ของเหลว (Liquid) และก๊าซ (Gas) หากเราพิจารณาน้ำแข็งซึ่งเป็นของแข็ง เมื่อได้รับพลังงานความร้อน พันธะระหว่างโมเลกุลจะอ่อนกำลังลงจนหลอมเหลวกลายเป็นน้ำในสถานะของเหลว และเมื่อให้พลังงานความร้อนต่อไปจนถึงจุดเดือด โมเลกุลของน้ำจะเคลื่อนที่ด้วยความเร็วสูงขึ้นจนหลุดออกจากกันและกลายเป็นไอน้ำในสถานะก๊าซ

คำถามสำคัญทางฟิสิกส์คือ หากเรายังคงป้อนพลังงานความร้อนหรือพลังงานไฟฟ้าเข้าไปยังก๊าซอย่างต่อเนื่องที่อุณหภูมิสูงขึ้นอย่างมหาศาลจะเกิดอะไรขึ้น พลังงานจลน์ที่สูงขึ้นจะทำให้อนุภาคในก๊าซชนกันอย่างรุนแรงจนทำให้อิเล็กตรอนวงนอกสุดหลุดออกจากอะตอมหรือโมเลกุล กระบวนการนี้เรียกว่า \textbf{การแตกตัวเป็นไอออน (Ionization)} ส่งผลให้กลุ่มก๊าซเดิมแปรสภาพกลายเป็นกลุ่มหมอกของอนุภาคที่มีประจุไฟฟ้า ประกอบด้วยอิเล็กตรอนอิสระ ไอออนบวก อะตอมหรือโมเลกุลที่ยังไม่แตกตัว ตลอดจนอนุมูลอิสระและรังสีแม่เหล็กไฟฟ้า สภาวะใหม่นี้เรียกว่า \textbf{พลาสมา (Plasma)} ซึ่งได้รับการยอมรับว่าเป็น \textbf{สถานะที่สี่ของสสาร (Fourth State of Matter)}

\begin{figure}[H]
\centering
\begin{tikzpicture}[scale=0.68, every node/.style={transform shape}]
    % Four states boxes
    \draw[rounded corners=6pt, fill=blue!10, draw=blue!60, line width=1pt] (-5.5, 0) rectangle (-3.0, 3.2);
    \node at (-4.25, 2.8) {\bfseries ของแข็ง};
    \node at (-4.25, 2.4) {\small (Solid)};
    \foreach \x in {-5.0, -4.5, -4.0, -3.5} {
        \foreach \y in {0.5, 1.0, 1.5, 2.0} {
            \fill[blue!70!black] (\x, \y) circle (3pt);
        }
    }

    \draw[->, >=Stealth, line width=1.5pt, red!70] (-2.8, 1.6) -- (-2.0, 1.6);
    \node[above, red!70] at (-2.4, 1.6) {\scriptsize +ความร้อน};

    \draw[rounded corners=6pt, fill=cyan!10, draw=cyan!60, line width=1pt] (-1.8, 0) rectangle (0.7, 3.2);
    \node at (-0.55, 2.8) {\bfseries ของเหลว};
    \node at (-0.55, 2.4) {\small (Liquid)};
    \foreach \pos in {(-1.4, 0.4), (-1.0, 0.7), (-0.5, 0.5), (0.2, 0.6), (-1.2, 1.2), (-0.7, 1.5), (0.1, 1.3), (-0.9, 2.0), (-0.2, 1.9)} {
        \fill[cyan!70!black] \pos circle (3.2pt);
    }

    \draw[->, >=Stealth, line width=1.5pt, red!70] (0.9, 1.6) -- (1.7, 1.6);
    \node[above, red!70] at (1.3, 1.6) {\scriptsize +ความร้อน};

    \draw[rounded corners=6pt, fill=yellow!10, draw=yellow!60!black, line width=1pt] (1.9, 0) rectangle (4.4, 3.2);
    \node at (3.15, 2.8) {\bfseries ก๊าซ};
    \node at (3.15, 2.4) {\small (Gas)};
    \foreach \pos in {(2.4, 0.8), (3.8, 0.6), (3.1, 1.4), (2.5, 2.1), (3.9, 1.9)} {
        \fill[yellow!60!black] \pos circle (3.5pt);
    }

    \draw[->, >=Stealth, line width=1.5pt, plasmaViolet] (4.6, 1.6) -- (5.4, 1.6);
    \node[above, plasmaViolet] at (5.0, 1.6) {\scriptsize +พลังงานไฟฟ้า};

    \draw[rounded corners=6pt, fill=plasmaViolet!10, draw=plasmaViolet, line width=1.2pt] (5.6, 0) rectangle (8.3, 3.2);
    \node at (6.95, 2.8) {\bfseries พลาสมา};
    \node at (6.95, 2.4) {\small (Plasma)};
    % Ions and electrons
    \foreach \pos in {(6.1, 0.8), (7.7, 0.7), (6.9, 1.5), (6.3, 2.1), (7.6, 2.0)} {
        \fill[red!70!black] \pos circle (4pt);
        \node[white] at \pos {\tiny +};
    }
    \foreach \pos in {(6.4, 0.5), (7.2, 0.9), (6.5, 1.2), (7.4, 1.6), (6.8, 2.2), (7.9, 2.4)} {
        \fill[blue!80!white] \pos circle (2pt);
        \node[white] at \pos {\scalebox{0.5}{$-$}};
    }
\end{tikzpicture}
\caption{การเปลี่ยนสถานะของสสารจากของแข็ง ของเหลว ก๊าซ สู่สถานะพลาสมาผ่านการแตกตัวเป็นไอออน}
\label{fig:four_states}
\end{figure}

แม้ว่าบนผิวโลกสสารส่วนใหญ่จะอยู่ในสถานะของแข็ง ของเหลว หรือก๊าซ แต่ในเอกภพโดยรวม สสารมากกว่าร้อยละ 99 อยู่ในสถานะพลาสมา เช่น ดวงอาทิตย์ ดวงดาวในกาแล็กซี เนบิวลา และสสารระหว่างดวงดาว บนผิวโลกเราสามารถพบพลาสมาตามธรรมชาติได้ในปรากฏการณ์ฟ้าผ่า (Lightning) และแสงเหนือแสงใต้ (Aurora)

\subsection{องค์ประกอบของพลาสมาความดันบรรยากาศ}

เมื่อก๊าซเปลี่ยนสถานะเป็นพลาสมา โดยเฉพาะพลาสมาที่เกิดจากการคายประจุในอากาศหรือก๊าซผสมที่ความดันบรรยากาศ ภายในพลาสมาจะมีกลุ่มอนุภาคและองค์ประกอบสำคัญ 6 ชนิด ได้แก่

\begin{enumerate}
    \item \textbf{อิเล็กตรอนอิสระ (Free Electrons หรือ $e^-$)} มีมวลน้อยมาก เคลื่อนที่ได้รวดเร็ว และเป็นตัวรับพลังงานหลักจากสนามไฟฟ้าภายนอกเพื่อถ่ายทอดพลังงานให้แก่อนุภาคอื่นๆ
    \item \textbf{ไอออนบวกและไอออนลบ (Positive and Negative Ions)} เกิดจากการที่อะตอมหรือโมเลกุลสูญเสียอิเล็กตรอน เช่น \ce{N2+}, \ce{O2+}, \ce{Ar+} หรือรับอิเล็กตรอนเพิ่ม เช่น \ce{O2-}, \ce{O-}
    \item \textbf{อะตอมและโมเลกุลในสถานะกระตุ้น (Excited Species and Metastables)} อนุภาคที่ได้รับพลังงานจากการชนแต่อิเล็กตรอนยังไม่หลุดออก เช่น \ce{O(^1D)}, \ce{N2(A^3{$\Sigma$}_u^+)} ซึ่งมีอายุขัยยาวนานและมีบทบาทสำคัญในการถ่ายโอนพลังงานเคมี
    \item \textbf{อนุมูลอิสระและสารว่องไวทางเคมี (Reactive Species หรือ RONS)} เกิดจากการแตกพันธะเคมีของก๊าซและไอน้ำในบรรยากาศ เช่น อนุมูลไฮดรอกซิล ($\cdot\ce{OH}$), โอโซน (\ce{O3}), ไฮโดรเจนเปอร์ออกไซด์ (\ce{H2O2}), ไนตริกออกไซด์ ($\cdot\ce{NO}$)
    \item \textbf{โฟตอนและรังสีแม่เหล็กไฟฟ้า (Photons and UV Radiation)} แสงที่ปลดปล่อยออกมาจากการเปลี่ยนระดับพลังงานของอิเล็กตรอน มีตั้งแต่แสงที่ตามองเห็นไปจนถึงรังสีอัลตราไวโอเลตสุญญากาศ (VUV)
    \item \textbf{สนามไฟฟ้าเฉพาะที่และคลื่นกระแทก (Local Electric Fields and Micro-shockwaves)} สนามไฟฟ้าความเข้มสูงระดับไมโครสเกลที่เกิดขึ้นระหว่างการเคลื่อนที่ของกลุ่มประจุ
\end{enumerate}

\section{นิยามและคุณสมบัติพื้นฐานทางฟิสิกส์ของพลาสมา}

ในทางฟิสิกส์ ก๊าซที่แตกตัวเป็นไอออนจะถือว่าเป็นพลาสมาได้อย่างแท้จริงก็ต่อเมื่อมีคุณสมบัติทางกายภาพครบ 3 ประการ ได้แก่ ความเป็นกลางทางไฟฟ้าโดยรวม การกำบังเดอบาย และพฤติกรรมรวมหมู่

\subsection{ความเป็นกลางทางไฟฟ้าโดยรวม}

แม้ว่าพลาสมาจะประกอบด้วยอนุภาคที่มีประจุไฟฟ้าลบและบวกจำนวนมหาศาล แต่เมื่อพิจารณาในปริมาตรมหภาค (Macroscopic scale) ประจุบวกและประจุลบจะหักล้างกันจนมีค่าสุทธิเป็นศูนย์ ความสัมพันธ์นี้เขียนได้ในรูป
\begin{equation}
    n_e \approx \sum_i Z_i n_i = n
\end{equation}
โดยที่ $n_e$ คือความหนาแน่นของอิเล็กตรอน ($m^{-3}$), $n_i$ คือความหนาแน่นของไอออนที่มีประจุ $Z_i$, และ $n$ เรียกว่าความหนาแน่นพลาสมา (Plasma Density)

\subsection{การกำบังเดอบายและความยาวเดอบาย}

เมื่อมีประจุไฟฟ้าแปลกปลอมหรือขั้วไฟฟ้าใส่เข้าไปในพลาสมา อนุภาคประจุตรงข้ามจะเคลื่อนที่เข้ามาล้อมรอบเพื่อบดบังสนามไฟฟ้าของประจุนั้น ปรากฏการณ์นี้เรียกว่า \textbf{การกำบังเดอบาย (Debye Shielding)} ระยะทางที่สนามไฟฟ้าลดลงเหลือ $1/e$ หรือประมาณร้อยละ 37 ของค่าเดิม เรียกว่า \textbf{ความยาวเดอบาย (Debye Length, $\lambda_D$)} ซึ่งมีนิยามตามสมการ
\begin{equation}
    \lambda_D = \sqrt{\frac{\varepsilon_0 k_B T_e}{n_e e^2}}
    \label{eq:debye_length}
\end{equation}
โดยที่ $\varepsilon_0 = 8.854 \times 10^{-12}\text{ F/m}$ คือสภาพยอมทางไฟฟ้าในสุญญากาศ, $k_B = 1.381 \times 10^{-23}\text{ J/K}$ คือค่าคงที่โบลต์ซมันน์, $T_e$ คืออุณหภูมิของอิเล็กตรอน (K), $n_e$ คือความหนาแน่นของอิเล็กตรอน ($m^{-3}$), และ $e = 1.602 \times 10^{-19}\text{ C}$ คือประจุของอิเล็กตรอน

\begin{figure}[H]
\centering
\begin{tikzpicture}[scale=0.95]
    % Central positive test charge
    \shade[inner color=red!80!white, outer color=red!90!black] (0,0) circle (0.35cm);
    \node[white, font=\bfseries] at (0,0) {$+Q$};

    % Debye sphere boundary
    \draw[dashed, line width=1.2pt, plasmaViolet!80] (0,0) circle (2.2cm);
    \draw[<->, >=Stealth, line width=1pt, plasmaNavy] (0,0) -- (1.55, 1.55);
    \node[above left, font=\small\bfseries, plasmaNavy] at (0.85, 0.85) {$\lambda_D$};

    % Surrounding electrons inside Debye sphere
    \foreach \angle in {25, 65, 110, 150, 195, 235, 280, 320, 350} {
        \fill[blue!70!white] (\angle:1.1cm) circle (2.5pt);
        \node[white] at (\angle:1.1cm) {\scalebox{0.5}{$-$}};
    }
    \foreach \angle in {15, 45, 85, 130, 175, 215, 260, 305, 340} {
        \fill[blue!70!white] (\angle:1.8cm) circle (2.5pt);
        \node[white] at (\angle:1.8cm) {\scalebox{0.5}{$-$}};
    }

    % Outer undisturbed quasi-neutral plasma
    \node[right, font=\small, plasmaNavy] at (2.4, 1.2) {บริเวณพลาสมาปกติ};
    \node[right, font=\footnotesize, gray!80!black] at (2.4, 0.7) {สนามไฟฟ้าภายนอกถูกบดบัง};
    \node[right, font=\footnotesize, plasmaViolet] at (2.4, -0.3) {ขอบเขตทรงกลมเดอบาย ($\lambda_D$)};
\end{tikzpicture}
\caption{การจัดเรียงกลุ่มหมอกอิเล็กตรอนล้อมรอบประจุทดสอบตามหลักการกำบังเดอบาย}
\label{fig:debye_shielding}
\end{figure}

ระบบก๊าซที่แตกตัวเป็นไอออนจะมีพฤติกรรมเป็นพลาสมาได้ต่อเมื่อขนาดทางกายภาพของระบบ ($L$) มีค่ามากกว่าความยาวเดอบายอย่างมีนัยสำคัญ นั่นคือ $L \gg \lambda_D$

\subsection{พฤติกรรมรวมหมู่และความถี่พลาสมา}

เนื่องจากแรงระหว่างอนุภาคที่มีประจุเป็นแรงคูลอมบ์ระยะไกล (Long-range Coulomb Force) การเคลื่อนที่ของอนุภาคหนึ่งจึงส่งผลต่ออนุภาคอื่นๆ พร้อมกัน เกิดเป็น \textbf{พฤติกรรมรวมหมู่ (Collective Behavior)} หากกลุ่มอิเล็กตรอนถูกรบกวนให้เคลื่อนที่ออกจากตำแหน่งสมดุล แรงดึงดูดทางไฟฟ้าจะดึงกลับเข้ามา เกิดการกวัดแกว่งรอบตำแหน่งสมดุลด้วย \textbf{ความถี่พลาสมาของอิเล็กตรอน (Electron Plasma Frequency, $\omega_{pe}$)}
\begin{equation}
    \omega_{pe} = \sqrt{\frac{n_e e^2}{\varepsilon_0 m_e}}
\end{equation}
โดยที่ $m_e = 9.109 \times 10^{-31}\text{ kg}$ คือมวลของอิเล็กตรอน

\subsection{ระดับการแตกตัวเป็นไอออนและอุณหภูมิในพลาสมา}

ระดับการแตกตัวเป็นไอออน ($\alpha$) นิยามเป็นสัดส่วนของความหนาแน่นไอออนเทียบกับอนุภาคทั้งหมด
\begin{equation}
    \alpha = \frac{n_i}{n_i + n_n}
\end{equation}
โดยที่ $n_n$ คือความหนาแน่นของอนุภาคที่เป็นกลาง ในพลาสมาความดันบรรยากาศอุณหภูมิต่ำ ค่า $\alpha$ มีค่าอยู่ในช่วง $10^{-7}$ ถึง $10^{-4}$ ซึ่งหมายความว่าโมเลกุลของอากาศส่วนใหญ่ยังคงเป็นก๊าซปกติ ทำให้อุณหภูมิรวมของก๊าซยังคงใกล้เคียงอุณหภูมิห้อง

ในทางฟิสิกส์พลาสมา นิยมวัดอุณหภูมิของอนุภาคในหน่วย \textbf{อิเล็กตรอนโวลต์ (Electron-Volt หรือ eV)} โดยแปลงเป็นเคลวินได้ตามความสัมพันธ์
\begin{equation}
    1\text{ eV} = \frac{e}{k_B} \approx 11,600\text{ K}
\end{equation}

\begin{focusbox}{ความสำคัญทางฟิสิกส์ของหน่วยอิเล็กตรอนโวลต์}
เมื่อกล่าวว่าพลาสมาเย็นมีอุณหภูมิอิเล็กตรอน $T_e = 2\text{ eV}$ หมายความว่าเฉพาะตัวอิเล็กตรอนมีพลังงานจลน์เฉลี่ยเทียบเท่าอุณหภูมิ $23,200\text{ K}$ ซึ่งมีพลังงานเพียงพอที่จะชนและตัดพันธะเคมีของก๊าซเพื่อสร้างอนุมูลอิสระ RONS ได้อย่างมหาศาล แต่เนื่องจากอิเล็กตรอนมีมวลน้อยมากและก๊าซส่วนใหญ่ไม่แตกตัว อุณหภูมิรวมที่เรารู้สึกเมื่อสัมผัสจึงยังคงอยู่ที่ $25 - 40^\circ\text{C}$
\end{focusbox}

\section{การจำแนกประเภทพลาสมาตามสมดุลทางอุณหพลศาสตร์}

การจำแนกพลาสมาตามความสัมพันธ์ระหว่างอุณหภูมิของอิเล็กตรอน ($T_e$), อุณหภูมิของไอออน ($T_i$), และอุณหภูมิของก๊าซ ($T_g$) แบ่งออกเป็น 2 ประเภทหลัก

\subsection{พลาสมาสมดุลความร้อน}

พลาสมาสมดุลความร้อน (Thermal Plasma) หรือพลาสมาร้อน เป็นสภาวะที่อนุภาคทุกชนิดมีพลังงานจลน์เฉลี่ยเท่ากันในระดับสมดุลอุณหพลศาสตร์เฉพาะที่ (Local Thermodynamic Equilibrium หรือ LTE)
\begin{equation}
    T_e \approx T_i \approx T_g \quad (5,000 - 50,000\text{ K})
\end{equation}
มีลักษณะเด่นคือให้พลังงานความร้อนสูงมหาศาล ความหนาแน่นอิเล็กตรอนสูง ($10^{21} - 10^{26}\text{ m}^{-3}$) เหมาะสำหรับงานเผาทำลายขยะติดเชื้อ การหลอมกากพิษให้กลายเป็นแก้วเสถียร (Vitrification) และการแก๊สซิฟิเคชันเพื่อเปลี่ยนชีวมวลเป็นพลังงาน

\subsection{พลาสมาไม่สมดุลความร้อนหรือพลาสมาเย็น}

พลาสมาไม่สมดุลความร้อน (Non-thermal Plasma) หรือพลาสมาเย็นความดันบรรยากาศ (Cold Atmospheric Plasma หรือ CAP) เป็นสภาวะที่พลังงานไฟฟ้าถูกถ่ายทอดให้แก่อิเล็กตรอนเป็นหลัก โดยที่ไอออนและโมเลกุลก๊าซยังคงมีพลังงานจลน์ต่ำ
\begin{equation}
    T_e \gg T_i \approx T_g \quad (T_e \approx 10,000 - 100,000\text{ K}, \text{ แต่ } T_g \approx 300 - 320\text{ K})
\end{equation}
ลักษณะเด่นคือไม่ทำลายเนื้อเยื่อพืชหรือวัสดุทางชีววิทยาที่บอบบาง แต่สามารถกระตุ้นปฏิกิริยาเคมีขั้นสูงได้ที่อุณหภูมิห้อง

\begin{table}[H]
\caption{การเปรียบเทียบคุณสมบัติระหว่างพลาสมาสมดุลความร้อนและพลาสมาไม่สมดุลความร้อน}
\label{tab:plasma_comparison}
\begin{tabularx}{\textwidth}{p{3.2cm}X X}
\toprule
\textbf{พารามิเตอร์} & \textbf{พลาสมาสมดุลความร้อน} & \textbf{พลาสมาเย็น (CAP)} \\
\midrule
สถานะอุณหพลศาสตร์ & สมดุลเฉพาะที่ ($T_e \approx T_i \approx T_g$) & ไม่สมดุล ($T_e \gg T_i \approx T_g$) \\
อุณหภูมิอิเล็กตรอน ($T_e$) & $1 - 5\text{ eV}$ ($10,000 - 50,000\text{ K}$) & $1 - 10\text{ eV}$ ($10,000 - 100,000\text{ K}$) \\
อุณหภูมิก๊าซรวม ($T_g$) & $5,000 - 50,000\text{ K}$ & $300 - 320\text{ K}$ ($25 - 45^\circ\text{C}$) \\
ความหนาแน่นอิเล็กตรอน & $10^{21} - 10^{26}\text{ m}^{-3}$ & $10^{14} - 10^{20}\text{ m}^{-3}$ \\
ระดับการแตกตัว ($\alpha$) & สูง ($10^{-1}$ ถึง $100\%$) & ต่ำมาก ($10^{-7}$ ถึง $10^{-4}$) \\
กลไกการทำงานหลัก & กระบวนการความร้อนสูง (Pyrolysis) & เคมีพลาสมาจากอนุภาค RONS \\
การประยุกต์ใช้ & กำจัดขยะอันตราย, หลอมกากโลหะ & กระตุ้นเมล็ด, ปุ๋ยน้ำ PAW, บำบัดมลพิษ \\
\bottomrule
\end{tabularx}
\end{table}

\section{ทำไมต้องใช้พลาสมาในงานเกษตรและสิ่งแวดล้อม}

การนำเทคโนโลยีพลาสมาเย็นความดันบรรยากาศมาใช้ในภาคเกษตรและสิ่งแวดล้อมมีจุดเด่นสำคัญ 5 ประการ

\begin{enumerate}
    \item \textbf{เทคโนโลยีสะอาดไร้สารตกค้าง} สารออกซิไดซ์และอนุมูลอิสระ RONS ที่เกิดขึ้นมีค่าครึ่งชีวิตสั้นและจะสลายตัวกลับคืนเป็นน้ำและออกซิเจนตามธรรมชาติ จึงไม่มีสารเคมีอันตรายตกค้างในดินหรือผลผลิต
    \item \textbf{การทำงานที่อุณหภูมิห้อง} สามารถใช้กับเนื้อเยื่อพืช เมล็ดพันธุ์ และผลผลิตสดหลังการเก็บเกี่ยวได้โดยไม่ทำลายคุณค่าทางอาหารและเซลล์พืช
    \item \textbf{กลไกหลายเป้าหมายลดปัญหาเชื้อดื้อยา} การรวมตัวกันของอนุมูล ROS, รังสี UV, สนามไฟฟ้า และไอออน จะทำลายเยื่อหุ้มเซลล์ ผนังเซลล์ และดีเอ็นเอของจุลินทรีย์พร้อมกัน ทำให้เชื้อก่อโรคไม่สามารถพัฒนาการดื้อยาได้เหมือนยาปฏิชีวนะเคมี
    \item \textbf{การเลียนแบบการตรึงไนโตรเจนตามธรรมชาติ} ใช้พลังงานไฟฟ้าเปลี่ยนไนโตรเจนและออกซิเจนในอากาศให้กลายเป็นสารประกอบไนเตรตและไนไตรต์ในน้ำ เกิดเป็น \textbf{น้ำกระตุ้นด้วยพลาสมา (Plasma-Activated Water หรือ PAW)} ซึ่งทำหน้าที่เป็นปุ๋ยไนโตรเจนธรรมชาติ
    \item \textbf{การกำจัดมลพิษย่อยสลายยากด้วยกระบวนการ AOPs} พลาสมาสร้างอนุมูล $\cdot\ce{OH}$ ซึ่งมีศักย์ไฟฟ้าออกซิเดชันสูงถึง $+2.80\text{ V}$ สามารถสลายโมเลกุลยาฆ่าแมลง สารอินทรีย์ระเหยง่าย และมลพิษกลุ่ม PFAS ได้อย่างสมบูรณ์
\end{enumerate}

\begin{agribox}{นวัตกรรมน้ำกระตุ้นด้วยพลาสมาเพื่อการเกษตรยั่งยืน}
การใช้น้ำกระตุ้นด้วยพลาสมา (PAW) เป็นการปฏิวัติการเกษตรด้วยการเปลี่ยนอากาศรอบตัวเป็นปุ๋ยไนโตรเจนและสารยับยั้งเชื้อราพร้อมใช้ โดยไม่ต้องพึ่งพากระบวนการฮาเบอร์-บอช (Haber-Bosch) ที่ปล่อยก๊าซคาร์บอนไดออกไซด์สูง เป็นหัวใจของเกษตรกรรมคาร์บอนต่ำในอนาคต
\end{agribox}

\clearpage

% =========================================================================
% คำถามทบทวนท้ายบทที่ 1 (ขึ้นหน้าใหม่อย่างสง่างาม)
% =========================================================================

\section*{คำถามทบทวนท้ายบทที่ 1}
\addcontentsline{toc}{section}{คำถามทบทวนท้ายบทที่ 1}

\begin{enumerate}
    \item จงอธิบายความแตกต่างระหว่างก๊าซในสภาวะปกติกับสสารในสถานะพลาสมา พร้อมทั้งระบุองค์ประกอบสำคัญ 6 ชนิดที่พบในพลาสมาความดันบรรยากาศ
    \item คำว่า Quasi-neutrality ในพลาสมามีความหมายทางกายภาพอย่างไร และเหตุใดขนาดทางกายภาพของระบบ ($L$) จึงต้องมีค่ามากกว่าความยาวเดอบาย ($\lambda_D$) เสมอ
    \item เหตุใดในพลาสมาเย็น (Cold Atmospheric Plasma) อิเล็กตรอนจึงมีอุณหภูมิสูงถึงหลายหมื่นเคลวิน ในขณะที่อุณหภูมิรวมของก๊าซยังคงอยู่ที่ระดับอุณหภูมิห้อง
    \item จงเปรียบเทียบข้อดีและข้อจำกัดของการใช้น้ำกระตุ้นด้วยพลาสมา (PAW) ในการตรึงไนโตรเจนเพื่อเป็นปุ๋ยพืช เมื่อเทียบกับการใช้ปุ๋ยเคมีไนโตรเจนสังเคราะห์
    \item จงวิเคราะห์ว่าหากต้องการกำจัดขยะติดเชื้อทางการแพทย์และหลอมกากเถ้าโลหะหนักไม่ให้รั่วไหล ควรเลือกใช้พลาสมาประเภทใด พร้อมอธิบายเหตุผลทางอุณหพลศาสตร์ประกอบ
\end{enumerate}

\clearpage
